% =========================================================
% Nested Interval Property
% =========================================================

\subsection{Nested Interval Property}


\begin{theorem}
\label{thm:nested-interval-property}
Let $(I_n)$ be a sequence of closed intervals in $\mathbb{R}$. If $(I_n)$ is a nested
interval sequence, then the intersection of the intervals is nonempty.
\end{theorem}

\begin{remark*}[Standard quantified statement]
\[
\begin{aligned}
&\forall (a_n) \subseteq \mathbb{R}\; \forall (b_n) \subseteq \mathbb{R}\;
\Biggl[
\Bigl(\forall n \in \mathbb{N}\; (a_n \le b_n)\Bigr)
\\
&\qquad\land
\Bigl(\forall n \in \mathbb{N}\; (I_n = [a_n,b_n])\Bigr)
\\
&\qquad\land
\Bigl(\forall n \in \mathbb{N}\; (I_{n+1} \subseteq I_n)\Bigr)
\\
&\qquad\rightarrow
\exists x \in \mathbb{R}\; \forall n \in \mathbb{N}\; (x \in I_n)
\Biggr].
\end{aligned}
\]
\end{remark*}

\begin{remark*}[Predicate statement]
\[
\forall (I_n)\;
\Bigl[
NestedIntervalSequence(I_n)
\rightarrow
NonemptySet\!\left(\bigcap_{n=1}^{\infty} I_n\right)
\Bigr].
\]
In this theorem, each $I_n$ is additionally assumed to be a closed interval in $\mathbb{R}$.
\end{remark*}

\begin{remark*}[Negated quantified statement]
\[
\begin{aligned}
&\exists (a_n) \subseteq \mathbb{R}\; \exists (b_n) \subseteq \mathbb{R}\;
\Biggl[
\Bigl(\forall n \in \mathbb{N}\; (a_n \le b_n)\Bigr)
\\
&\qquad\land
\Bigl(\forall n \in \mathbb{N}\; (I_n = [a_n,b_n])\Bigr)
\\
&\qquad\land
\Bigl(\forall n \in \mathbb{N}\; (I_{n+1} \subseteq I_n)\Bigr)
\\
&\qquad\land
\forall x \in \mathbb{R}\; \exists n \in \mathbb{N}\; (x \notin I_n)
\Biggr].
\end{aligned}
\]
This means there exists a nested sequence of closed intervals in $\mathbb{R}$ whose total
intersection is empty. Such a failure would indicate a defect in the ambient system, not in
the nesting construction itself.
\end{remark*}

\input{volume-iii/analysis/bounding/notes/completeness/figure-nested-intervals}

\begin{remark*}[Interpretation]
The Nested Interval Property says that a decreasing sequence of closed intervals in
$\mathbb{R}$ cannot shrink in such a way that every point is eventually lost. If each interval
lies inside the previous one, then at least one real number remains in all of them.

The important point is that nesting prevents the intervals from wandering apart, while
closedness prevents the limiting point from dropping out. The theorem says that this process
of repeated containment still leaves a common point in the end.
\end{remark*}